\documentclass[aspectratio=169]{beamer}

\usetheme{Boadilla}
\usecolortheme{default}
\setbeamertemplate{navigation symbols}{}
\setbeamertemplate{footline}[frame number]

\usepackage{amsmath, amssymb, booktabs}

\newcommand{\Ztwo}{\mathbb{Z}_2}
\newcommand{\Ftwo}{\mathbb{F}_2}

\title[2-adic Numbers in Cryptography]{2-adic Numbers in Cryptography}
\subtitle{From carries to stream ciphers}
\author{Siddharth K}
\date{May 2026}

\begin{document}

% 1
\begin{frame}
  \titlepage
\end{frame}

% 2
\begin{frame}{The main idea}
\large
A long bit string can be treated as one 2-adic number:
\[
  a_0,a_1,a_2,\ldots
  \quad\longleftrightarrow\quad
  \alpha = a_0+a_1 2+a_2 2^2+\cdots \in \Ztwo.
\]

\vspace{0.4cm}
Today I will explain:
\begin{enumerate}
  \item why this is natural for cryptography,
  \item what a feedback with carry shift register is,
  \item why its output is a rational 2-adic number.
\end{enumerate}
\end{frame}

% 3
\begin{frame}{Why cryptography needs bit streams}
Alice has the message written as bits:
\[
  P_0, P_1, P_2, \ldots, \quad P_i \in \{0, 1\}.
\]

A stream cipher uses a keystream:
\[
  K_0,K_1,K_2,\ldots
\]

Alice sends
\[
  C_i = P_i + K_i \pmod 2.
\]

Bob decrypts by doing the same operation:
\[
  P_i = C_i + K_i \pmod 2.
\]

\begin{block}{The problem}
Alice and Bob need the same keystream, but Eve should not be able to predict it.
\end{block}
\end{frame}

% 4
\begin{frame}{2-adic integers}
A 2-adic integer is an infinite binary sum
\[
  a_0+a_1 2+a_2 2^2+a_3 2^3+\cdots,
  \qquad a_i\in\{0,1\}.
\]

\vspace{0.2cm}
The key difference from real numbers:
\[
  2^n \to 0 \quad \text{in the 2-adic metric.}
\]

So infinite sums of powers of $2$ can converge in $\Ztwo$.
\end{frame}

% 5
\begin{frame}{Two examples}

\vspace{0.2cm}
First example:
\[
  111111\cdots = 1+2+2^2+2^3+\cdots = -1 \quad \text{in }\Ztwo.
\]

Second example:
\[
  101010\cdots = 1+2^2+2^4+\cdots
\]
\[
  = 1+4+4^2+\cdots
  = \frac{1}{1-4}
  = -\frac{1}{3}.
\]

\begin{block}{Takeaway}
Repeating bit strings often become rational 2-adic numbers.
\end{block}
\end{frame}

% 6
\begin{frame}{Ordinary shift registers}
A linear feedback shift register, or LFSR, creates new bits using a linear rule.

Example:
\[
  a_{n+3}=a_{n+1}+a_n \pmod 2.
\]

\vspace{0.3cm}
Why people like LFSRs:
\begin{itemize}
  \item they are fast,
  \item they are simple,
  \item they can make long periodic sequences.
\end{itemize}

\begin{alertblock}{Problem}
They are linear over $\Ftwo$. With enough output, an attacker may recover the rule.
\end{alertblock}
\end{frame}

% 7
\begin{frame}{Feedback with carry shift registers}
An FCSR is like an LFSR, but it remembers carries.

Choose tap values $q_1,\ldots,q_r\in\{0,1\}$ and a memory value $m_{n-1}$.
At step $n$, compute
\[
  \sigma_n=q_1a_{n-1}+q_2a_{n-2}+\cdots+q_ra_{n-r}+m_{n-1}.
\]

Then set
\[
  a_n \equiv \sigma_n \pmod 2,
  \qquad
  m_n=\frac{\sigma_n-a_n}{2}.
\]

\begin{block}{Simple picture}
LFSR: $1+1=0$ mod $2$, and the extra $1$ disappears.\newline
FCSR: $1+1=10_2$, so the output is $0$ and the carry is $1$.
\end{block}
\end{frame}

% 8
\begin{frame}{The connection integer}
The taps define a number called the connection integer:
\[
  q=-1+q_1 2+q_2 2^2+\cdots+q_r2^r.
\]

Usually we take $q>0$, so $q$ is odd.

\vspace{0.4cm}
\begin{block}{Main theorem}
Let
\[
  \alpha=a_0+a_1 2+a_2 2^2+\cdots
\]
be the 2-adic number made from the output of an FCSR with connection integer $q$.
Then there is an integer $p$ such that
\[
  \alpha=\frac{p}{q} \quad \text{in }\Ztwo.
\]
\end{block}
\end{frame}

% 9
\begin{frame}{Proof idea}
Start with
\[
  \alpha=x+\sum_{n=r}^{\infty}a_n2^n,
  \qquad
  x=a_0+a_1 2+\cdots+a_{r-1}2^{r-1}.
\]

From the FCSR rule,
\[
  a_n=q_1a_{n-1}+\cdots+q_ra_{n-r}+m_{n-1}-2m_n.
\]

Substitute this into the tail:
\[
\alpha=x+
\sum_{i=1}^r q_i\sum_{n=r}^{\infty}a_{n-i}2^n
+
\sum_{n=r}^{\infty}m_{n-1}2^n
-
\sum_{n=r}^{\infty}m_n2^{n+1}.
\]
\end{frame}

% 10
\begin{frame}{Proof idea}
For the bit sums, re-index:
\[
  \sum_{n=r}^{\infty}a_{n-i}2^n
  =2^i\left(\alpha-X_i\right),
\]
where
\[
  X_i=\sum_{j=0}^{r-i-1}a_j2^j
\]
is just a finite starting piece.

\vspace{0.2cm}
For the memory sums, take finite sums first:
\[
  \sum_{n=r}^{N}m_{n-1}2^n-
  \sum_{n=r}^{N}m_n2^{n+1}
  =m_{r-1}2^r-m_N2^{N+1}.
\]

As $N\to\infty$, the last term goes to $0$ in $\Ztwo$.
\end{frame}

% 11
\begin{frame}{Proof idea}
Putting the pieces together gives
\[
  \alpha=x+(q+1)\alpha-
  \sum_{i=1}^{r}q_i2^iX_i+2^rm_{r-1}.
\]

Move the $\alpha$ terms to one side:
\[
  q\alpha=-x+
  \sum_{i=1}^{r}q_i2^iX_i-2^rm_{r-1}.
\]

The right side is an ordinary integer. Call it $p$.

So
\[
  q\alpha=p,
  \qquad \text{and therefore}\qquad
  \alpha=\frac{p}{q}.
\]

\begin{block}{What happened?}
The infinite tail collapsed to finite data: the starting bits and the carry.
\end{block}
\end{frame}

% 12
\begin{frame}{Example: connection integer $q=5$}
If $q=5$, then
\[
  q+1=6=2+4,
\]
so the rule is
\[
  \sigma_n=a_{n-1}+a_{n-2}+m_{n-1}.
\]

Start with
\[
  a_0=1,\qquad a_1=1,\qquad m_1=0.
\]

\begin{center}
\begin{tabular}{c|c|c|c}
$n$ & $\sigma_n$ & new bit $a_n$ & new memory $m_n$ \\
\hline
2 & $1+1+0=2$ & $0$ & $1$ \\
3 & $0+1+1=2$ & $0$ & $1$ \\
4 & $0+0+1=1$ & $1$ & $0$ \\
5 & $1+0+0=1$ & $1$ & $0$
\end{tabular}
\end{center}

So the sequence starts
\[
  11001100\cdots.
\]
\end{frame}

% 13
\begin{frame}{Example: its 2-adic value}
The sequence
\[
  11001100\cdots
\]
means
\[
  \alpha=1+2+2^4+2^5+2^8+2^9+\cdots.
\]

Group the repeating blocks:
\[
  \alpha=(1+2)(1+2^4+2^8+\cdots).
\]

Since $2^4=16$,
\[
  \alpha=3(1+16+16^2+\cdots)
  =3\cdot \frac{1}{1-16}
  =-\frac{1}{5}.
\]

This matches the theorem: the connection integer was $5$.
\end{frame}

% 14
\begin{frame}{What this means for cryptography}
FCSRs can generate keystreams, but they are not magic.

\vspace{0.2cm}
The theorem says the output has structure:
\[
  \alpha=\frac{p}{q}.
\]

If Eve sees the first $N$ bits, she knows
\[
  \alpha \pmod{2^N}.
\]

A 2-adic rational approximation attack asks:
\[
  \text{Can we find small }p,q\text{ with }\frac{p}{q}\equiv \alpha \pmod{2^N}?
\]

\begin{alertblock}{Security issue}
The carry makes the generator less linear, but the 2-adic structure can still guide attacks.
\end{alertblock}
\end{frame}

% 15
\begin{frame}{Takeaways}
\begin{enumerate}
  \item A bit string can be viewed as a 2-adic integer.
  \item FCSRs are shift registers that remember carries.
  \item The output of an FCSR is a rational 2-adic number $p/q$.
  \item This gives both a construction idea and an attack idea.
\end{enumerate}

\vfill
\centering
Thanks for listening!
\end{frame}

\end{document}